\documentclass{article}

%\usepackage{corl_2026}
\usepackage[final]{corl_2026}

\usepackage{booktabs}
\usepackage{amsmath}
\usepackage{graphicx}
\usepackage{subcaption}
\usepackage{url}

\title{SoccerTwos: PPO, Behavioral Cloning, and Reward Shaping}

\author{
  Shaoyu Zeng\\
  Georgia Institute of Technology
  \And
  Haoran Liu\\
  Georgia Institute of Technology
}

\begin{document}
\maketitle

\noindent\textbf{Code repository:}
\url{https://github.com/Carbonadooo/soccer-twos-starter}

\section{Methodology}

We trained three agents for SoccerTwos, a sparse-reward 2-vs-2 soccer task with
$+1$ for scoring and $-1$ for conceding.

\textbf{Agent 1: PPO vs.\ random opponent.}
We used PPO~\citep{schulman2017proximal} in Ray RLlib~\citep{liang2018rllib}.
One policy jointly controls both teammates in \texttt{team\_vs\_policy} mode
against a random opponent.
The training script used for this agent was
\texttt{train\_ray\_team\_vs\_random.py}.
PPO optimizes the clipped objective
\[
L^{\text{CLIP}}(\theta)=
\mathbb{E}\left[\min(r_t(\theta)\hat A_t,
\text{clip}(r_t(\theta),1-\epsilon,1+\epsilon)\hat A_t)\right].
\]
We chose this as the baseline because it is the simplest end-to-end RL setup.

\textbf{Agent 2: Behavioral Cloning (BC).}
This agent uses supervised imitation learning instead of RL.
The training script used for this stage was \texttt{distill\_bc\_obs.py}.
We collected 300 episodes of expert state-action pairs from the teacher baseline,
normalized each observation dimension using z-score normalization,
\[
\tilde s_d=\frac{s_d-\mu_d}{\sigma_d+10^{-8}},
\]
and trained a two-layer MLP with three action heads using cross-entropy loss.
Our hypothesis was that imitation would bypass the sparse-reward exploration
problem and produce a stronger starting policy than random initialization.

\textbf{Agent 3: BC + PPO + reward shaping.}
We initialized PPO from BC weights, then fine-tuned against the teacher baseline.
The training script used for this agent was
\texttt{train\_ray\_bc\_finetune\_vs\_baseline\_shaped.py}.
The main reward modification was dense shaping added to the sparse goal reward.
The shaped reward used the actual positions and velocities returned by the
environment info dictionary.
First, we predicted a short-horizon ball position and rewarded the increase in
normalized forward ball progress, which acts as a potential-style shaping term.
Second, the teammate nearest to the predicted ball position received a small
ball-contest bonus based on an exponential distance decay, encouraging one player
to actively pressure or control the ball.
Third, when the ball was in our own half, the team received a defensive coverage
bonus if at least one player stayed between the ball and our goal along the ball
lane, and a small penalty otherwise.
Fourth, the support player was rewarded for maintaining useful width and depth in
attack, or for covering the goal lane in defense.
Finally, an anti-crowding penalty discouraged both teammates from collapsing to
the same point, while a lateral-separation bonus encouraged more useful spacing.
In short, the reward modification tried to turn soccer tactics that matter for
winning into dense step-by-step learning signals:
\begin{itemize}
\item ball-progress potential,
\item nearest-player ball contest reward,
\item defensive coverage bonus,
\item support-positioning bonus,
\item anti-crowding and spacing terms.
\end{itemize}
The motivation was that these terms should provide denser learning signal and
speed up convergence relative to plain PPO.

\begin{table}[h]
\centering
\caption{Final training settings used for the RL agents.}
\label{tab:hparams}
\begin{tabular}{lcc}
\toprule
\textbf{Parameter} & \textbf{Agent 1} & \textbf{Agent 3} \\
\midrule
Algorithm & PPO & PPO \\
Library & RLlib 1.4 & RLlib 1.4 \\
Opponent & Random & Teacher baseline \\
Workers & 8 & 8 \\
Envs/worker & 3 & 2 \\
Rollout fragment & 500 & 500 \\
Train batch size & 12000 & 8000 \\
Learning rate & default & $2\times10^{-5}$ \\
Clip parameter & 0.2 & 0.1 \\
Entropy coefficient & 0.001 & default \\
Hidden layers & [512, 512] & [512, 512] \\
Total timesteps & 3.6M & 7.2M \\
\bottomrule
\end{tabular}
\end{table}

\section{Experimental Results}

Figure~\ref{fig:curves} shows the reward-vs.-timesteps curves for the two RL
agents.
Agent~2 is not shown there because it is trained with supervised learning rather
than RL, so reward-vs.-steps is not the right metric for that stage.

\begin{figure}[h]
  \centering
  \begin{subfigure}{0.48\linewidth}
    \centering
    \includegraphics[width=\linewidth]{figures/agent1_reward_curves.png}
    \caption{Agent 1}
  \end{subfigure}
  \hfill
  \begin{subfigure}{0.48\linewidth}
    \centering
    \includegraphics[width=\linewidth]{figures/agent3_reward_curves.png}
    \caption{Agent 3}
  \end{subfigure}
  \caption{Training curves. X-axis: total environment timesteps. Y-axis: episode
  reward mean.}
  \label{fig:curves}
\end{figure}

\begin{table}[h]
\centering
\caption{Evaluation against the teacher baseline over 200 episodes.}
\label{tab:eval}
\begin{tabular}{lcc}
\toprule
\textbf{Agent} & \textbf{Wins / 200} & \textbf{Win Rate} \\
\midrule
Agent 1: PPO vs.\ random & 31 / 200 & 0.155 \\
Agent 2: BC only & 107 / 200 & 0.535 \\
Agent 3: BC + PPO + shaping & 136 / 200 & 0.680 \\
\bottomrule
\end{tabular}
\end{table}

\section{Analysis and Discussion}

The results support our original hypothesis.
Plain PPO against a random opponent learned useful behavior, but it transferred
poorly to the stronger teacher baseline, winning only 15.5\% of games.
This suggests that sparse rewards and weak opponents were enough to learn basic
control, but not enough to produce a strong competitive policy.

Behavioral cloning performed much better than Agent~1, reaching a 53.5\% win
rate against the teacher baseline.
This supports the idea that imitation learning is an effective way to overcome
the exploration problem in sparse-reward environments.

Agent~3 performed best at 68.0\% win rate.
This indicates that our modification was useful: BC provided a strong
initialization, and reward shaping gave PPO denser feedback than the raw
goal-only reward.
Compared with Agent~2, Agent~3 shows that RL fine-tuning on top of imitation can
improve beyond simply copying the baseline.

The reward curves also suggest a convergence effect.
Agent~1 improves steadily, while Agent~3 reaches a higher-performing regime
faster but also shows some instability later in training.
That pattern is consistent with the stronger supervision from reward shaping:
faster learning, but also greater sensitivity to checkpoint choice.

Some of this behavior is also connected to our hyperparameter choices.
For Agent~1, we kept the setup fairly close to a standard PPO baseline, but we
did increase the clip parameter to 0.2 and added a small entropy coefficient of
0.001.
This likely encouraged broader exploration against the random opponent and helped
the policy keep improving steadily instead of collapsing too early.
For Agent~3, we changed more of the defaults: the learning rate was reduced to
$2\times10^{-5}$, the clip parameter was tightened to 0.1, the number of vector
environments per worker was reduced from 3 to 2, and the train batch size was
reduced from 12000 to 8000.
These choices were meant to make fine-tuning from BC more conservative, so PPO
would refine the imitation policy instead of overwriting it too aggressively in
the harder teacher-baseline setting.
The final results are consistent with that design: Agent~3 beats the teacher
baseline much more often than Agent~1, but its reward curve is less smooth, which
suggests the stronger shaping signal and more delicate fine-tuning regime made
checkpoint selection more important.

\clearpage
\bibliography{report}

\end{document}
